\documentclass[11pt]{article}

\usepackage[margin=1in]{geometry}
\usepackage{amsmath,amssymb}
\usepackage{booktabs}
\usepackage{graphicx}
\usepackage{subcaption}
\usepackage{xcolor}
\usepackage{array}
\usepackage{multirow}
\usepackage{longtable}
\usepackage{pdflscape}
\usepackage{hyperref}
\usepackage{float}
\usepackage{placeins}
\usepackage{adjustbox}
\usepackage{enumitem}

\hypersetup{colorlinks=true, linkcolor=blue, citecolor=blue, urlcolor=blue}

\title{Report0520: Step1b--Step1c Update\\
\large From Perturbed FY Generalization to Matched SPO+ Surrogate Ablation}
\author{Weikang}
\date{May 20, 2026}

% -----------------------------------------------------------------------------
% Path conventions
% -----------------------------------------------------------------------------
% This report is designed to be compiled from the repository root.
% If you move this .tex file into surrogate_experiment_results/ and compile there,
% change the following line to: \newcommand{\RepoRoot}{..}
\newcommand{\RepoRoot}{.}

\newcommand{\StepOneBPlotRoot}{\RepoRoot/surrogate_experiment_results/Step1b/plot_results}
\newcommand{\StepOneBCPlotRoot}{\RepoRoot/surrogate_experiment_results/plot_results_step1bc}
\newcommand{\StepOneBValRoot}{\RepoRoot/surrogate_experiment_results/Step1b/remote_results/formal_M16_2stage500_e2e500_s10_val2000}
\newcommand{\StepOneCRoot}{\RepoRoot/surrogate_experiment_results/Step1c/remote_results/formal_spoplus_ablation_val2000}
\newcommand{\StepOneCTrain}[1]{\StepOneCRoot/train_size=#1}

\newcommand{\maybeincludegraphics}[2][]{%
  \IfFileExists{#2}{%
    \includegraphics[#1]{#2}%
  }{%
    \fbox{\begin{minipage}{0.92\linewidth}\centering
    Missing figure file:\\[0.4em]
    \texttt{\detokenize{#2}}
    \end{minipage}}%
  }%
}

\newcommand{\artifactpath}[1]{\path{#1}}
\newcommand{\spoplus}{SPO\textsuperscript{+}}
\newcommand{\fy}{Fenchel--Young}
\newcommand{\gap}{synthetic-label decision gap}
\newcommand{\normgap}{normalized decision gap}

\begin{document}
\maketitle

\section{Purpose of Report0520}

Report0513 established the Step1a--Step1b sequence. Step1a was an in-sample
landscape diagnostic for a two-dimensional reward probe, and Step1b upgraded
that diagnostic into a held-out generalization study comparing a two-stage MSE
baseline with perturbed \fy{} decision-focused training.

This report documents the work completed after Report0513. The main new
scientific step is \textbf{Step1c}, a matched surrogate-ablation experiment:
we keep the Step1b data protocol, model class, checkpointing logic, and test
metrics fixed, but replace the end-to-end \fy{} surrogate with a reward-max
\spoplus{} surrogate. The report is designed to be presentation-friendly: the
main text explains what changed, how to read each figure, and what the new data
support or do not yet support.

\paragraph{Key question for this week.}
\begin{quote}
Compared with the Report0513 Step1b setup, does replacing perturbed \fy{} with
reward-max \spoplus{} change held-out KEP decision quality under the same
synthetic noisy-linear benchmark protocol?
\end{quote}

\paragraph{Metric language.}
As in Report0513, all labels are synthetic rewards. The main metric is therefore
called \emph{synthetic-label decision gap}, \emph{oracle objective gap}, or
\emph{normalized decision gap}; it should not be described as clinical regret.

\section{Summary of New Work Since Report0513}

Table~\ref{tab:weekly-change-summary} gives the high-level changes. The rest of
this report walks through the experiment design, tables, and figures in the same
order.

\begin{table}[H]
\centering
\caption{Report0520 changes relative to Report0513.}
\label{tab:weekly-change-summary}
\resizebox{\textwidth}{!}{%
\begin{tabular}{p{0.23\linewidth}p{0.36\linewidth}p{0.36\linewidth}}
\toprule
Component & Report0513 state & New state in Report0520 \\
\midrule
Validation protocol & Step1b used the original 400-graph validation split for the formal archive discussed in Report0513. & Step1b and Step1c now use a larger validation override, \artifactpath{dataset/processed/step1_noisy_linear_sigma010_validation2000_seed20260519}, for stricter checkpoint comparison. \\
\midrule
Main end-to-end surrogate & Perturbed \fy{} only. & Added reward-max \spoplus{} as Step1c, while preserving Step1b's dataset, split, two-feature model class, and evaluation metrics. \\
\midrule
Train sizes & Step1b completed mostly \(n=50,200,600\); \(n=1200\) was still a planned run in the report text. & Step1c has formal \spoplus{} runs for \(n=50,200,600,1200\). The strict Step1b validation2000 archive currently has \(n=50,200,600\) locally available; Step1b \(n=1200\) is still absent from the strict combined archive. \\
\midrule
Unseen evaluation & Smaller unseen noisy-linear evaluation. & Added a larger unseen noisy-linear evaluation archive. \\
\midrule
Diagnostics & Aggregate Step1b figures and tables. & Added post-hoc Step1b dashboards, reward calibration plots, trajectory dashboards, selection-suboptimality analysis, per-graph diagnostics, and combined Step1b/Step1c comparison figures. \\
\midrule
Engineering reliability & Step1b training/evaluation scripts. & Added Step1c \spoplus{} trainer, one-pass large-unseen evaluator, final comparison plotting, and unit tests for \spoplus{}, post-hoc diagnostics, unseen evaluation, and combined plotting. \\
\bottomrule
\end{tabular}}
\end{table}

\section{Experimental Protocol}

\subsection{Benchmark and model class}

The benchmark remains the controlled noisy-linear synthetic KEP setting used in
Report0513:
\[
    w^{\mathrm{syn}}_e
    =
    \max\left\{0,\,(10u_e + 5c_e)(1+\xi_e)\right\},
    \qquad \sigma=0.10,
\]
where \(u_e\) is edge utility and \(c_e\) is recipient cPRA. All compared methods
use the same two-feature linear reward probe:
\[
    \hat w_e(\theta)=\theta_1 u_e + \theta_2 c_e .
\]
Thus the comparison is not about model capacity. It is about how different
training losses move the same model class through the downstream KEP decision
landscape.

\subsection{Datasets and archive status}

\begin{table}[H]
\centering
\caption{Datasets and result archives used for Report0520.}
\label{tab:dataset-archive-status}
\resizebox{\textwidth}{!}{%
\begin{tabular}{p{0.23\linewidth}p{0.39\linewidth}p{0.31\linewidth}}
\toprule
Item & Path / archive & Current status \\
\midrule
Training source & \artifactpath{dataset/processed/step1_noisy_linear_sigma010} & Same benchmark as Report0513. \\
Validation override & \artifactpath{dataset/processed/step1_noisy_linear_sigma010_validation2000_seed20260519} & New default validation set for strict Step1b/Step1c checkpoint selection. \\
Held-out test & Original 400-graph test split from \artifactpath{results/step1b_splits/master_split_seed=42.json} & Used for the main held-out comparison. \\
Large unseen test & Internal processed large-unseen archive & New large out-of-split noisy-linear evaluation. \\
Step1b strict FY archive & \artifactpath{surrogate_experiment_results/Step1b/remote_results/formal_M16_2stage500_e2e500_s10_val2000} & Available for \(n=50,200,600\) in the current combined report. \\
Step1c SPO+ archive & \artifactpath{surrogate_experiment_results/Step1c/remote_results/formal_spoplus_ablation_val2000} & Available for \(n=50,200,600,1200\). \\
Combined output & \artifactpath{surrogate_experiment_results/plot_results_step1bc} & Contains final combined figures and summary tables. \\
\bottomrule
\end{tabular}}
\end{table}

\subsection{Methods compared}

\begin{table}[H]
\centering
\caption{Model checkpoints compared in Report0520.}
\label{tab:methods-compared}
\resizebox{\textwidth}{!}{%
\begin{tabular}{p{0.26\linewidth}p{0.36\linewidth}p{0.29\linewidth}}
\toprule
Method label in figures & Training objective & Checkpoint selection rule \\
\midrule
2stage (val MSE) & Fit synthetic edge rewards by MSE, then optimize KEP using predicted weights. & Lowest validation MSE. \\
FY (val gap) & Perturbed \fy{} decision-focused training. & Lowest validation synthetic-label decision gap. Diagnostic / oracle-style selection. \\
FY (val FY) & Perturbed \fy{} decision-focused training. & Lowest validation \fy{} surrogate objective. Main FY surrogate-selection checkpoint. \\
SPO+ (val gap) & Reward-max \spoplus{} decision-focused training. & Lowest validation synthetic-label decision gap. Diagnostic / oracle-style selection. \\
SPO+ (val SPO+) & Reward-max \spoplus{} decision-focused training. & Lowest validation \spoplus{} objective. Main SPO+ surrogate-selection checkpoint. \\
\bottomrule
\end{tabular}}
\end{table}

\section{Step1c: Reward-Max SPO+ Surrogate}

\subsection{Reward-max formulation}

The original SPO literature is often written for cost minimization. KEP in this
project is reward maximization:
\[
    y^\star(w) \in \arg\max_{y\in\mathcal{Y}} w^\top y.
\]
For Step1c, signs are converted by treating costs as \(c=-w\). For one graph,
let \(w\) be synthetic label rewards, \(X\) be the two-feature edge matrix,
\(\hat w=X\theta\), and \(y^\star=y^\star(w)\). The reward-max \spoplus{} loss is
\[
    L_{\mathrm{SPO+}}(\theta)
    =
    \max_{y\in\mathcal{Y}} (2\hat w - w)^\top y
    -2\hat w^\top y^\star
    +w^\top y^\star .
\]
The shifted objective is
\[
    w^{\mathrm{adv}}=2\hat w-w,
    \qquad
    y^{\mathrm{adv}}=y^\star(w^{\mathrm{adv}}),
\]
and the subgradient is
\[
    \nabla_\theta L_{\mathrm{SPO+}}(\theta)
    =2X^\top\left(y^{\mathrm{adv}}-y^\star\right).
\]

\paragraph{Interpretation.}
The \fy{} update compares the oracle solution with an average perturbed
optimizer solution. The \spoplus{} update compares the oracle solution with a
single deterministic adversarial shifted-objective solution. Thus Step1c asks
whether this deterministic margin-style challenger gives a better training and
checkpoint-selection signal than perturbed \fy{}.

\subsection{Step1c diagnostics recorded}

In addition to the final test gaps, the \spoplus{} runs record raw/normalized
\spoplus{} loss, decision gap, \(\|\theta\|\), and two solution-equality rates:
\[
    \Pr\left[y^\star(2\hat w-w)=y^\star(w)\right],
    \qquad
    \Pr\left[y^\star(\hat w)=y^\star(w)\right].
\]
These are useful during discussion because they separate two effects: whether
SPO+'s adversarial challenger has been defeated, and whether the current
prediction already induces the oracle KEP solution.

\section{Main Combined Figures}

The final combined figures are generated under
\artifactpath{surrogate_experiment_results/plot_results_step1bc}. They compare
2stage, FY, and SPO+ under the matched validation2000 protocol whenever the
corresponding archive is available.

\subsection{Mean performance overview}

\begin{figure}[H]
\centering
\maybeincludegraphics[width=0.98\textwidth]{\StepOneBCPlotRoot/figure0_mean_normalized_gap_no_error_bars.png}
\caption{Mean normalized decision gap without uncertainty bars. Lower is better.
This figure is useful as the first visual summary because it shows the trend by
train size for the five selected checkpoints.}
\label{fig:step1bc-figure0}
\end{figure}

\paragraph{How to read Figure~\ref{fig:step1bc-figure0}.}
The left panel is the held-out 400-graph test split. The right panel is the new
large unseen noisy-linear test. The key visual pattern is that \spoplus{}
selected by validation \spoplus{} loss has a strong held-out400 signal, while
the large-unseen panel is much closer to a tie across methods.

\begin{figure}[H]
\centering
\maybeincludegraphics[width=0.98\textwidth]{\StepOneBCPlotRoot/figure1_mean_normalized_gap_heldout400_large_unseen.png}
\caption{Mean normalized decision gap with per-graph bootstrap 95\% confidence
intervals. Lower is better.}
\label{fig:step1bc-figure1}
\end{figure}

\paragraph{How to read Figure~\ref{fig:step1bc-figure1}.}
This is the more statistically cautious version of
Figure~\ref{fig:step1bc-figure0}. On heldout400, the \spoplus{} validation-loss
checkpoint is consistently lower than the 2stage baseline. On the large unseen test, the
bars are much closer together, so the safer statement is that \spoplus{} is
competitive but not yet a robust unseen-data improvement.

\subsection{Per-graph distribution}

\begin{figure}[H]
\centering
\maybeincludegraphics[width=0.98\textwidth]{\StepOneBCPlotRoot/figure2_per_graph_gap_boxplots_heldout400_large_unseen.png}
\caption{Per-graph normalized decision-gap distributions for heldout400 and the
large unseen test. Lower is better.}
\label{fig:step1bc-figure2}
\end{figure}

\paragraph{How to read Figure~\ref{fig:step1bc-figure2}.}
Many KEP graphs have zero or near-zero decision gap for several methods. The
boxplots therefore help answer whether an average improvement is broad across
many graphs or driven by a small subset of graphs where the induced KEP solution
changes.

\subsection{Checkpoint and parameter diagnostics}

\begin{figure}[H]
\centering
\begin{subfigure}{0.49\textwidth}
\centering
\maybeincludegraphics[width=\linewidth]{\StepOneBCPlotRoot/figure3_checkpoint_epochs.png}
\caption{Selected epochs.}
\end{subfigure}
\hfill
\begin{subfigure}{0.49\textwidth}
\centering
\maybeincludegraphics[width=\linewidth]{\StepOneBCPlotRoot/figure4_selected_theta_endpoints.png}
\caption{Selected parameter endpoints.}
\end{subfigure}
\caption{Checkpoint-selection diagnostics. These figures explain which points on
each training trajectory are used in the final comparison.}
\label{fig:step1bc-figures3-4}
\end{figure}

\paragraph{How to read Figure~\ref{fig:step1bc-figures3-4}.}
The epoch plot shows whether each surrogate's validation loss selects early or
late checkpoints. The parameter plot connects back to Step1a: MSE stays near the
clean coefficient \([10,5]\), while decision-focused methods can select smaller
or larger parameter vectors depending on the surrogate and validation criterion.

\begin{figure}[H]
\centering
\maybeincludegraphics[width=0.98\textwidth]{\StepOneBCPlotRoot/figure5_loss_curves.png}
\caption{Training and validation loss curves for 2stage reward fitting, FY, and
SPO+. Dashed curves are train losses and solid curves are validation losses.}
\label{fig:step1bc-figure5}
\end{figure}

\paragraph{How to read Figure~\ref{fig:step1bc-figure5}.}
The three panels use different loss scales, so they should not be compared by
vertical magnitude. Instead, use them to check whether each training objective is
stable and whether the validation curve gives a plausible checkpoint-selection
signal.

\FloatBarrier
\section{Numerical Results: Held-out 400-Graph Test}

\subsection{Strict Step1b FY rerun under validation2000}

Table~\ref{tab:step1b-heldout400-val2000} summarizes the strict Step1b FY
baseline currently available under the validation2000 protocol. Compared with
Report0513, this is the fairer FY baseline for Step1c because both use the same
larger validation set for checkpoint selection. The local combined archive
currently contains \(n=50,200,600\), not \(n=1200\), for this strict FY rerun.

\begin{table}[H]
\centering
\caption{Step1b validation2000 FY rerun on heldout400. Lower gap is better;
positive paired improvement means the candidate improves over 2stage on the
same test graphs.}
\label{tab:step1b-heldout400-val2000}
\resizebox{\textwidth}{!}{%
\begin{tabular}{rllrrrrrr}
\toprule
Train size & Method & Selected by & Epoch & \(\theta_1\) & \(\theta_2\) & Mean gap & Norm. gap & Paired improvement \\
\midrule
50 & 2stage & val MSE & 500 & 9.8707 & 5.1130 & 0.8480 & 0.006110 & -- \\
50 & FY & val gap & 50 & 5.4807 & 2.8132 & 0.8220 & 0.005918 & 0.0260 \\
50 & FY & val FY & 440 & 11.8304 & 5.0208 & 0.7587 & 0.005467 & 0.0894 \\
\midrule
200 & 2stage & val MSE & 500 & 9.8848 & 5.1038 & 0.8451 & 0.006095 & -- \\
200 & FY & val gap & 240 & 10.3559 & 5.2299 & 0.8073 & 0.005809 & 0.0378 \\
200 & FY & val FY & 500 & 12.2682 & 6.0649 & 0.7730 & 0.005629 & 0.0721 \\
\midrule
600 & 2stage & val MSE & 500 & 9.8899 & 5.0963 & 0.8223 & 0.005909 & -- \\
600 & FY & val gap & 130 & 8.3220 & 4.2029 & 0.8073 & 0.005809 & 0.0150 \\
600 & FY & val FY & 490 & 12.2328 & 5.8586 & 0.7507 & 0.005497 & 0.0716 \\
\bottomrule
\end{tabular}}
\end{table}

\paragraph{Reading Table~\ref{tab:step1b-heldout400-val2000}.}
The validation2000 rerun preserves the qualitative finding from Report0513: the
FY checkpoint selected by validation FY loss is better than the 2stage baseline
on the original heldout400 split for the completed train sizes. This motivates
using FY as the strict baseline when comparing SPO+.

\subsection{Step1c SPO+ heldout400 results}

\begin{table}[H]
\centering
\caption{Step1c SPO+ results on heldout400. Lower gap is better; positive paired
improvement means the candidate improves over 2stage on the same test graphs.}
\label{tab:step1c-heldout400}
\resizebox{\textwidth}{!}{%
\begin{tabular}{rllrrrrrrl}
\toprule
Train size & Method & Selected by & Epoch & \(\theta_1\) & \(\theta_2\) & Mean gap & Norm. gap & Paired improvement & 95\% CI \\
\midrule
50 & 2stage & val MSE & 500 & 9.8707 & 5.1130 & 0.8480 & 0.006110 & -- & -- \\
50 & SPO+ & val gap & 100 & 6.7245 & 3.4149 & 0.8073 & 0.005809 & 0.0407 & [-0.0001, 0.1000] \\
50 & SPO+ & val SPO+ & 480 & 9.1155 & 4.1380 & 0.7401 & 0.005442 & 0.1079 & [0.0347, 0.1960] \\
\midrule
200 & 2stage & val MSE & 500 & 9.8848 & 5.1038 & 0.8451 & 0.006095 & -- & -- \\
200 & SPO+ & val gap & 130 & 7.0334 & 3.5855 & 0.8185 & 0.005891 & 0.0266 & [-0.0076, 0.0830] \\
200 & SPO+ & val SPO+ & 500 & 9.0462 & 4.4396 & 0.7706 & 0.005610 & 0.0745 & [0.0170, 0.1527] \\
\midrule
600 & 2stage & val MSE & 500 & 9.8899 & 5.0963 & 0.8223 & 0.005909 & -- & -- \\
600 & SPO+ & val gap & 110 & 6.8532 & 3.5499 & 0.8480 & 0.006110 & -0.0258 & [-0.0734, 0.0000] \\
600 & SPO+ & val SPO+ & 500 & 9.1956 & 4.5117 & 0.7706 & 0.005610 & 0.0517 & [0.0044, 0.1112] \\
\midrule
1200 & 2stage & val MSE & 500 & 9.8848 & 5.0949 & 0.8441 & 0.006085 & -- & -- \\
1200 & SPO+ & val gap & 130 & 7.0369 & 3.6820 & 0.8450 & 0.006093 & -0.0009 & [-0.0098, 0.0082] \\
1200 & SPO+ & val SPO+ & 500 & 9.0009 & 4.4856 & 0.7846 & 0.005697 & 0.0594 & [0.0033, 0.1312] \\
\bottomrule
\end{tabular}}
\end{table}

\paragraph{Reading Table~\ref{tab:step1c-heldout400}.}
The main Step1c signal is the row selected by validation SPO+ loss. It improves
over the 2stage baseline for all four train sizes on heldout400. The diagnostic
row selected by validation decision gap is less stable: it helps at \(n=50\) and
\(n=200\), but is weak or worse at larger train sizes. Therefore the cleanest
Step1c message is that \emph{validation SPO+ loss is a useful surrogate-selection
criterion on the original held-out split}.

\section{Numerical Results: Large Unseen Noisy-Linear Test}

The large unseen evaluation is the new out-of-split stress test added after
Report0513. It is larger than the previous unseen evaluation and is therefore a
stricter check of whether the heldout400 improvements generalize to newly
generated KEP graphs from the same label regime.

\begin{table}[H]
\centering
\caption{Step1c SPO+ results on the large unseen noisy-linear test. Lower
gap is better. Positive paired improvement means the candidate improves over
2stage on the same unseen graphs.}
\label{tab:step1c-large-unseen}
\resizebox{\textwidth}{!}{%
\begin{tabular}{rllrrrrrrl}
\toprule
Train size & Method & Selected by & Epoch & \(\theta_1\) & \(\theta_2\) & Mean gap & Norm. gap & Paired improvement & 95\% CI \\
\midrule
50 & 2stage & val MSE & 500 & 9.8707 & 5.1130 & 0.6790 & 0.005118 & -- & -- \\
50 & SPO+ & val gap & 100 & 6.7245 & 3.4149 & 0.6801 & 0.005127 & -0.0010 & [-0.0044, 0.0025] \\
50 & SPO+ & val SPO+ & 480 & 9.1155 & 4.1380 & 0.6909 & 0.005220 & -0.0119 & [-0.0240, 0.0003] \\
\midrule
200 & 2stage & val MSE & 500 & 9.8848 & 5.1038 & 0.6799 & 0.005125 & -- & -- \\
200 & SPO+ & val gap & 130 & 7.0334 & 3.5855 & 0.6799 & 0.005126 & 0.0000 & [-0.0021, 0.0021] \\
200 & SPO+ & val SPO+ & 500 & 9.0462 & 4.4396 & 0.6794 & 0.005135 & 0.0005 & [-0.0066, 0.0078] \\
\midrule
600 & 2stage & val MSE & 500 & 9.8899 & 5.0963 & 0.6797 & 0.005124 & -- & -- \\
600 & SPO+ & val gap & 110 & 6.8532 & 3.5499 & 0.6787 & 0.005116 & 0.0010 & [-0.0011, 0.0034] \\
600 & SPO+ & val SPO+ & 500 & 9.1956 & 4.5117 & 0.6795 & 0.005136 & 0.0001 & [-0.0069, 0.0074] \\
\midrule
1200 & 2stage & val MSE & 500 & 9.8848 & 5.0949 & 0.6797 & 0.005124 & -- & -- \\
1200 & SPO+ & val gap & 130 & 7.0369 & 3.6820 & 0.6802 & 0.005129 & -0.0005 & [-0.0051, 0.0036] \\
1200 & SPO+ & val SPO+ & 500 & 9.0009 & 4.4856 & 0.6769 & 0.005109 & 0.0028 & [-0.0028, 0.0084] \\
\bottomrule
\end{tabular}}
\end{table}

\paragraph{Reading Table~\ref{tab:step1c-large-unseen}.}
The large unseen evaluation is much more cautious than heldout400. SPO+ is close
to 2stage, but the improvements are small and the confidence intervals still
cross zero. Thus the report should not claim robust SPO+ dominance on newly
generated unseen graphs. The better statement is: SPO+ gives a strong heldout400
signal and remains competitive on the large unseen test, but the unseen advantage is not
established yet.

\FloatBarrier
\section{Step1b Diagnostic Figures Added Since Report0513}

This section collects the additional Step1b diagnostic figures generated after
Report0513. They are useful when walking through the week with a supervisor:
first show the combined result figures in Section~4, then use the following
plots to explain why the aggregate bars behave the way they do.

\subsection{Aggregate Step1b result plots}

\begin{figure}[H]
\centering
\maybeincludegraphics[width=\textwidth]{\StepOneBPlotRoot/01_normalized_decision_gap_by_dataset.png}
\caption{Step1b mean normalized decision gap by dataset. Lower is better.}
\label{fig:step1b-01}
\end{figure}

\begin{figure}[H]
\centering
\maybeincludegraphics[width=\textwidth]{\StepOneBPlotRoot/02_raw_decision_gap_by_dataset.png}
\caption{Step1b raw synthetic-label decision gap by dataset. Lower is better.}
\label{fig:step1b-02}
\end{figure}

\begin{figure}[H]
\centering
\maybeincludegraphics[width=\textwidth]{\StepOneBPlotRoot/03_achieved_oracle_ratio_by_dataset.png}
\caption{Step1b achieved/oracle objective ratio by dataset. Higher is better.}
\label{fig:step1b-03}
\end{figure}

\begin{figure}[H]
\centering
\maybeincludegraphics[width=\textwidth]{\StepOneBPlotRoot/04_paired_improvement_over_2stage.png}
\caption{Step1b direct method-performance comparison. Lower mean normalized gap
is better; paired statistics remain in the summary tables.}
\label{fig:step1b-04}
\end{figure}

\paragraph{How to read Figures~\ref{fig:step1b-01}--\ref{fig:step1b-04}.}
These figures are the Step1b-only view. They explain the FY baseline before
adding SPO+. They show that FY selected by validation FY loss is promising on
heldout400, but the advantage weakens on unseen data and shifted stress tests.

\subsection{Step1b parameter and training diagnostics}

\begin{figure}[H]
\centering
\maybeincludegraphics[width=0.82\textwidth]{\StepOneBPlotRoot/05_selected_parameter_endpoints.png}
\caption{Step1b selected parameter endpoints in the two-dimensional linear-probe
space.}
\label{fig:step1b-05}
\end{figure}

\begin{figure}[H]
\centering
\maybeincludegraphics[width=\textwidth]{\StepOneBPlotRoot/06_parameter_trajectories.png}
\caption{Step1b parameter trajectories across train sizes.}
\label{fig:step1b-06}
\end{figure}

\begin{figure}[H]
\centering
\maybeincludegraphics[width=\textwidth]{\StepOneBPlotRoot/07_08_09_training_curves_combined.png}
\caption{Step1b combined training curves: 2stage MSE, FY objective, and decision
gap.}
\label{fig:step1b-07-08-09}
\end{figure}

\begin{figure}[H]
\centering
\maybeincludegraphics[width=\textwidth]{\StepOneBPlotRoot/07_2stage_mse_loss_curves.png}
\caption{Step1b 2stage MSE loss curves.}
\label{fig:step1b-07}
\end{figure}

\begin{figure}[H]
\centering
\maybeincludegraphics[width=\textwidth]{\StepOneBPlotRoot/08_e2e_fy_objective_curves.png}
\caption{Step1b FY objective curves.}
\label{fig:step1b-08}
\end{figure}

\begin{figure}[H]
\centering
\maybeincludegraphics[width=\textwidth]{\StepOneBPlotRoot/09_e2e_decision_gap_curves.png}
\caption{Step1b e2e decision-gap curves.}
\label{fig:step1b-09}
\end{figure}

\begin{figure}[H]
\centering
\maybeincludegraphics[width=\textwidth]{\StepOneBPlotRoot/10_e2e_validation_tradeoff.png}
\caption{Step1b validation FY objective versus validation decision gap.}
\label{fig:step1b-10}
\end{figure}

\paragraph{How to read Figures~\ref{fig:step1b-05}--\ref{fig:step1b-10}.}
These plots connect back to Step1a. MSE remains close to the clean coefficient
\([10,5]\), while FY can select different parameter scales. The validation
tradeoff plot is especially important: it asks whether the surrogate objective
is aligned enough with validation decision quality to be used for checkpoint
selection.

\subsection{Step1b per-graph, calibration, and post-hoc diagnostics}

\begin{figure}[H]
\centering
\begin{subfigure}{0.32\textwidth}
\centering
\maybeincludegraphics[width=\linewidth]{\StepOneBPlotRoot/11_heldout400_gap_boxplots.png}
\caption{Heldout400.}
\end{subfigure}
\hfill
\begin{subfigure}{0.32\textwidth}
\centering
\maybeincludegraphics[width=\linewidth]{\StepOneBPlotRoot/12_unseen1000_gap_boxplots.png}
\caption{Unseen1000.}
\end{subfigure}
\hfill
\begin{subfigure}{0.32\textwidth}
\centering
\maybeincludegraphics[width=\linewidth]{\StepOneBPlotRoot/13_realistic2000_gap_boxplots.png}
\caption{Realistic2000.}
\end{subfigure}
\caption{Step1b per-graph normalized decision-gap boxplots.}
\label{fig:step1b-11-13}
\end{figure}

\begin{figure}[H]
\centering
\begin{subfigure}{0.32\textwidth}
\centering
\maybeincludegraphics[width=\linewidth]{\StepOneBPlotRoot/14_heldout400_reward_calibration.png}
\caption{Heldout400.}
\end{subfigure}
\hfill
\begin{subfigure}{0.32\textwidth}
\centering
\maybeincludegraphics[width=\linewidth]{\StepOneBPlotRoot/15_unseen1000_reward_calibration.png}
\caption{Unseen1000.}
\end{subfigure}
\hfill
\begin{subfigure}{0.32\textwidth}
\centering
\maybeincludegraphics[width=\linewidth]{\StepOneBPlotRoot/16_realistic2000_reward_calibration.png}
\caption{Realistic2000.}
\end{subfigure}
\caption{Step1b reward calibration plots. Edges are binned by predicted reward;
the diagonal is perfect calibration against synthetic labels.}
\label{fig:step1b-14-16}
\end{figure}

\paragraph{How to read Figures~\ref{fig:step1b-11-13}--\ref{fig:step1b-14-16}.}
The boxplots show whether average gap differences come from many graphs or a
small subset of difficult graphs. The calibration plots separate reward-fitting
quality from decision quality: a model can be reward-calibrated yet not always
select the best KEP solution under the synthetic-label objective.

\begin{figure}[H]
\centering
\begin{subfigure}{0.49\textwidth}
\centering
\maybeincludegraphics[width=\linewidth]{\StepOneBPlotRoot/heldout400_epoch_diagnostics.png}
\caption{Heldout400 epoch diagnostics.}
\end{subfigure}
\hfill
\begin{subfigure}{0.49\textwidth}
\centering
\maybeincludegraphics[width=\linewidth]{\StepOneBPlotRoot/unseen1000_epoch_diagnostics.png}
\caption{Unseen1000 epoch diagnostics.}
\end{subfigure}

\vspace{0.8em}
\begin{subfigure}{0.72\textwidth}
\centering
\maybeincludegraphics[width=\linewidth]{\StepOneBPlotRoot/realistic2000_epoch_diagnostics.png}
\caption{Realistic2000 epoch diagnostics.}
\end{subfigure}
\caption{Step1b epoch-wise diagnostics on test datasets.}
\label{fig:step1b-epoch-diagnostics}
\end{figure}

\begin{figure}[H]
\centering
\begin{subfigure}{0.49\textwidth}
\centering
\maybeincludegraphics[width=\linewidth]{\StepOneBPlotRoot/17_trajectory_dashboard_train_size_50.png}
\caption{Train size 50.}
\end{subfigure}
\hfill
\begin{subfigure}{0.49\textwidth}
\centering
\maybeincludegraphics[width=\linewidth]{\StepOneBPlotRoot/17_trajectory_dashboard_train_size_200.png}
\caption{Train size 200.}
\end{subfigure}

\vspace{0.8em}
\begin{subfigure}{0.49\textwidth}
\centering
\maybeincludegraphics[width=\linewidth]{\StepOneBPlotRoot/17_trajectory_dashboard_train_size_600.png}
\caption{Train size 600.}
\end{subfigure}
\hfill
\begin{subfigure}{0.49\textwidth}
\centering
\maybeincludegraphics[width=\linewidth]{\StepOneBPlotRoot/17_trajectory_dashboard_train_size_1200.png}
\caption{Train size 1200.}
\end{subfigure}
\caption{Step1b post-hoc trajectory dashboards. Vertical markers indicate
validation-selected checkpoints and post-hoc best evaluated epochs on different
datasets.}
\label{fig:step1b-17-dashboard}
\end{figure}

\begin{figure}[H]
\centering
\maybeincludegraphics[width=0.88\textwidth]{\StepOneBPlotRoot/18_selection_suboptimality.png}
\caption{Step1b checkpoint-selection suboptimality. This is a post-hoc diagnostic:
it measures the gap between a selected checkpoint and the best evaluated epoch
on a given dataset.}
\label{fig:step1b-18-selection-suboptimality}
\end{figure}

\begin{figure}[H]
\centering
\begin{subfigure}{0.32\textwidth}
\centering
\maybeincludegraphics[width=\linewidth]{\StepOneBPlotRoot/19_heldout400_per_graph_diagnostics.png}
\caption{Heldout400.}
\end{subfigure}
\hfill
\begin{subfigure}{0.32\textwidth}
\centering
\maybeincludegraphics[width=\linewidth]{\StepOneBPlotRoot/19_unseen1000_per_graph_diagnostics.png}
\caption{Unseen1000.}
\end{subfigure}
\hfill
\begin{subfigure}{0.32\textwidth}
\centering
\maybeincludegraphics[width=\linewidth]{\StepOneBPlotRoot/19_realistic2000_per_graph_diagnostics.png}
\caption{Realistic2000.}
\end{subfigure}
\caption{Step1b per-graph post-hoc diagnostics.}
\label{fig:step1b-19-per-graph}
\end{figure}

\paragraph{How to read Figures~\ref{fig:step1b-epoch-diagnostics}--\ref{fig:step1b-19-per-graph}.}
These figures are diagnostic rather than primary evidence. The trajectory
dashboards help explain checkpoint timing. The selection-suboptimality plot shows
whether validation criteria choose a point close to the best evaluated epoch.
The per-graph diagnostics focus on nonzero-gap cases, because many graphs induce
identical KEP decisions across methods.

\FloatBarrier
\section{Step1c Individual Run Diagnostics}

The final comparison figures aggregate across train sizes, but it is useful to
also keep the individual Step1c training curves for each run. Each train size has
a 2stage MSE curve and a SPO+ surrogate curve under
\artifactpath{surrogate_experiment_results/Step1c/remote_results/formal_spoplus_ablation_val2000/train_size=<n>/plots}.

\begin{figure}[H]
\centering
\begin{subfigure}{0.49\textwidth}
\centering
\maybeincludegraphics[width=\linewidth]{\StepOneCTrain{50}/plots/2stage_mse_loss.png}
\caption{2stage MSE, \(n=50\).}
\end{subfigure}
\hfill
\begin{subfigure}{0.49\textwidth}
\centering
\maybeincludegraphics[width=\linewidth]{\StepOneCTrain{50}/plots/spoplus_loss.png}
\caption{SPO+ loss, \(n=50\).}
\end{subfigure}
\caption{Step1c training curves for train size 50.}
\label{fig:step1c-curves-50}
\end{figure}

\begin{figure}[H]
\centering
\begin{subfigure}{0.49\textwidth}
\centering
\maybeincludegraphics[width=\linewidth]{\StepOneCTrain{200}/plots/2stage_mse_loss.png}
\caption{2stage MSE, \(n=200\).}
\end{subfigure}
\hfill
\begin{subfigure}{0.49\textwidth}
\centering
\maybeincludegraphics[width=\linewidth]{\StepOneCTrain{200}/plots/spoplus_loss.png}
\caption{SPO+ loss, \(n=200\).}
\end{subfigure}
\caption{Step1c training curves for train size 200.}
\label{fig:step1c-curves-200}
\end{figure}

\begin{figure}[H]
\centering
\begin{subfigure}{0.49\textwidth}
\centering
\maybeincludegraphics[width=\linewidth]{\StepOneCTrain{600}/plots/2stage_mse_loss.png}
\caption{2stage MSE, \(n=600\).}
\end{subfigure}
\hfill
\begin{subfigure}{0.49\textwidth}
\centering
\maybeincludegraphics[width=\linewidth]{\StepOneCTrain{600}/plots/spoplus_loss.png}
\caption{SPO+ loss, \(n=600\).}
\end{subfigure}
\caption{Step1c training curves for train size 600.}
\label{fig:step1c-curves-600}
\end{figure}

\begin{figure}[H]
\centering
\begin{subfigure}{0.49\textwidth}
\centering
\maybeincludegraphics[width=\linewidth]{\StepOneCTrain{1200}/plots/2stage_mse_loss.png}
\caption{2stage MSE, \(n=1200\).}
\end{subfigure}
\hfill
\begin{subfigure}{0.49\textwidth}
\centering
\maybeincludegraphics[width=\linewidth]{\StepOneCTrain{1200}/plots/spoplus_loss.png}
\caption{SPO+ loss, \(n=1200\).}
\end{subfigure}
\caption{Step1c training curves for train size 1200.}
\label{fig:step1c-curves-1200}
\end{figure}

\paragraph{How to read Figures~\ref{fig:step1c-curves-50}--\ref{fig:step1c-curves-1200}.}
These are sanity checks for the individual Step1c runs. They show whether the
MSE and SPO+ training objectives were numerically stable before interpreting the
final heldout/unseen metrics.

\FloatBarrier
\section{Complete Artifact Roadmap}

Table~\ref{tab:artifact-roadmap} lists the figure and table artifacts that should
be kept together with this report. The table is intentionally comprehensive so
that during a meeting we can jump from a high-level claim to the corresponding
figure or CSV file.

\begin{longtable}{p{0.26\linewidth}p{0.49\linewidth}p{0.20\linewidth}}
\caption{Report0520 artifact roadmap.}\label{tab:artifact-roadmap}\\
\toprule
Group & Artifact path & Role in report \\
\midrule
\endfirsthead
\toprule
Group & Artifact path & Role in report \\
\midrule
\endhead
\midrule
\multicolumn{3}{r}{Continued on next page}\\
\endfoot
\bottomrule
\endlastfoot
Combined final comparison & \artifactpath{surrogate_experiment_results/plot_results_step1bc/figure0_mean_normalized_gap_no_error_bars.png} / .pdf & Main mean trend without uncertainty bars. \\
Combined final comparison & \artifactpath{surrogate_experiment_results/plot_results_step1bc/figure1_mean_normalized_gap_heldout400_large_unseen.png} / .pdf & Main grouped-bar comparison with 95\% CI. \\
Combined final comparison & \artifactpath{surrogate_experiment_results/plot_results_step1bc/figure2_per_graph_gap_boxplots_heldout400_large_unseen.png} / .pdf & Per-graph distribution. \\
Combined final comparison & \artifactpath{surrogate_experiment_results/plot_results_step1bc/figure3_checkpoint_epochs.png} / .pdf & Selected epoch diagnostic. \\
Combined final comparison & \artifactpath{surrogate_experiment_results/plot_results_step1bc/figure4_selected_theta_endpoints.png} / .pdf & Selected parameter diagnostic. \\
Combined final comparison & \artifactpath{surrogate_experiment_results/plot_results_step1bc/figure5_loss_curves.png} / .pdf & Training/validation loss overview. \\
Combined final comparison & \artifactpath{surrogate_experiment_results/plot_results_step1bc/combined_step1bc_summary.csv} & Source table for combined summary rows. \\
Combined final comparison & \artifactpath{surrogate_experiment_results/plot_results_step1bc/combined_step1bc_per_graph.csv} & Source table for per-graph rows; may be gitignored because it is large. \\
\midrule
Step1b aggregate figures & \artifactpath{surrogate_experiment_results/Step1b/plot_results/01_normalized_decision_gap_by_dataset.png} / .pdf & Step1b normalized gap. \\
Step1b aggregate figures & \artifactpath{surrogate_experiment_results/Step1b/plot_results/02_raw_decision_gap_by_dataset.png} / .pdf & Step1b raw gap. \\
Step1b aggregate figures & \artifactpath{surrogate_experiment_results/Step1b/plot_results/03_achieved_oracle_ratio_by_dataset.png} / .pdf & Step1b achieved/oracle ratio. \\
Step1b aggregate figures & \artifactpath{surrogate_experiment_results/Step1b/plot_results/04_paired_improvement_over_2stage.png} / .pdf & Step1b direct method performance. \\
Step1b diagnostics & \artifactpath{surrogate_experiment_results/Step1b/plot_results/05_selected_parameter_endpoints.png} / .pdf & Parameter endpoints. \\
Step1b diagnostics & \artifactpath{surrogate_experiment_results/Step1b/plot_results/06_parameter_trajectories.png} / .pdf & Parameter trajectories. \\
Step1b diagnostics & \artifactpath{surrogate_experiment_results/Step1b/plot_results/07_08_09_training_curves_combined.png} / .pdf & Combined training curves. \\
Step1b diagnostics & \artifactpath{surrogate_experiment_results/Step1b/plot_results/07_2stage_mse_loss_curves.png} / .pdf & MSE loss curves. \\
Step1b diagnostics & \artifactpath{surrogate_experiment_results/Step1b/plot_results/08_e2e_fy_objective_curves.png} / .pdf & FY objective curves. \\
Step1b diagnostics & \artifactpath{surrogate_experiment_results/Step1b/plot_results/09_e2e_decision_gap_curves.png} / .pdf & FY decision-gap curves. \\
Step1b diagnostics & \artifactpath{surrogate_experiment_results/Step1b/plot_results/10_e2e_validation_tradeoff.png} / .pdf & Validation FY vs decision-gap tradeoff. \\
Step1b per-graph & \artifactpath{surrogate_experiment_results/Step1b/plot_results/11_heldout400_gap_boxplots.png} / .pdf & Heldout400 boxplot. \\
Step1b per-graph & \artifactpath{surrogate_experiment_results/Step1b/plot_results/12_unseen1000_gap_boxplots.png} / .pdf & Unseen1000 boxplot. \\
Step1b per-graph & \artifactpath{surrogate_experiment_results/Step1b/plot_results/13_realistic2000_gap_boxplots.png} / .pdf & Realistic2000 boxplot. \\
Step1b calibration & \artifactpath{surrogate_experiment_results/Step1b/plot_results/14_heldout400_reward_calibration.png} / .pdf & Heldout reward calibration. \\
Step1b calibration & \artifactpath{surrogate_experiment_results/Step1b/plot_results/15_unseen1000_reward_calibration.png} / .pdf & Unseen reward calibration. \\
Step1b calibration & \artifactpath{surrogate_experiment_results/Step1b/plot_results/16_realistic2000_reward_calibration.png} / .pdf & Realistic reward calibration. \\
Step1b epoch diagnostics & \artifactpath{surrogate_experiment_results/Step1b/plot_results/heldout400_epoch_diagnostics.png} / .pdf & Epoch diagnostics for heldout400. \\
Step1b epoch diagnostics & \artifactpath{surrogate_experiment_results/Step1b/plot_results/unseen1000_epoch_diagnostics.png} / .pdf & Epoch diagnostics for unseen1000. \\
Step1b epoch diagnostics & \artifactpath{surrogate_experiment_results/Step1b/plot_results/realistic2000_epoch_diagnostics.png} / .pdf & Epoch diagnostics for realistic2000. \\
Step1b post-hoc & \artifactpath{surrogate_experiment_results/Step1b/plot_results/17_trajectory_dashboard_train_size_50.png} / .pdf & Trajectory dashboard, \(n=50\). \\
Step1b post-hoc & \artifactpath{surrogate_experiment_results/Step1b/plot_results/17_trajectory_dashboard_train_size_200.png} / .pdf & Trajectory dashboard, \(n=200\). \\
Step1b post-hoc & \artifactpath{surrogate_experiment_results/Step1b/plot_results/17_trajectory_dashboard_train_size_600.png} / .pdf & Trajectory dashboard, \(n=600\). \\
Step1b post-hoc & \artifactpath{surrogate_experiment_results/Step1b/plot_results/17_trajectory_dashboard_train_size_1200.png} / .pdf & Trajectory dashboard, \(n=1200\). \\
Step1b post-hoc & \artifactpath{surrogate_experiment_results/Step1b/plot_results/18_selection_suboptimality.png} / .pdf and \artifactpath{18_selection_suboptimality.csv} & Selection-suboptimality diagnostic. \\
Step1b post-hoc & \artifactpath{surrogate_experiment_results/Step1b/plot_results/19_heldout400_per_graph_diagnostics.png} / .pdf & Heldout per-graph diagnostics. \\
Step1b post-hoc & \artifactpath{surrogate_experiment_results/Step1b/plot_results/19_unseen1000_per_graph_diagnostics.png} / .pdf & Unseen per-graph diagnostics. \\
Step1b post-hoc & \artifactpath{surrogate_experiment_results/Step1b/plot_results/19_realistic2000_per_graph_diagnostics.png} / .pdf & Realistic per-graph diagnostics. \\
Step1b post-hoc & \artifactpath{surrogate_experiment_results/Step1b/plot_results/19_per_graph_diagnostics_summary.csv} & Summary table for per-graph diagnostics. \\
\midrule
Step1c train-size 50 & \artifactpath{surrogate_experiment_results/Step1c/remote_results/formal_spoplus_ablation_val2000/train_size=50/metrics/test_summary.csv} & Heldout400 SPO+ summary. \\
Step1c train-size 50 & Internal large-unseen summary CSV & Large-unseen SPO+ summary. \\
Step1c train-size 50 & \artifactpath{surrogate_experiment_results/Step1c/remote_results/formal_spoplus_ablation_val2000/train_size=50/metrics/spoplus_loss_curve.csv} & SPO+ training curve table. \\
Step1c train-size 50 & \artifactpath{surrogate_experiment_results/Step1c/remote_results/formal_spoplus_ablation_val2000/train_size=50/plots/2stage_mse_loss.png}; \artifactpath{plots/spoplus_loss.png} & Individual run curves. \\
Step1c train-size 200 & \artifactpath{surrogate_experiment_results/Step1c/remote_results/formal_spoplus_ablation_val2000/train_size=200/metrics/test_summary.csv} & Heldout400 SPO+ summary. \\
Step1c train-size 200 & Internal large-unseen summary CSV & Large-unseen SPO+ summary. \\
Step1c train-size 200 & \artifactpath{surrogate_experiment_results/Step1c/remote_results/formal_spoplus_ablation_val2000/train_size=200/metrics/spoplus_loss_curve.csv} & SPO+ training curve table. \\
Step1c train-size 200 & \artifactpath{surrogate_experiment_results/Step1c/remote_results/formal_spoplus_ablation_val2000/train_size=200/plots/2stage_mse_loss.png}; \artifactpath{plots/spoplus_loss.png} & Individual run curves. \\
Step1c train-size 600 & \artifactpath{surrogate_experiment_results/Step1c/remote_results/formal_spoplus_ablation_val2000/train_size=600/metrics/test_summary.csv} & Heldout400 SPO+ summary. \\
Step1c train-size 600 & Internal large-unseen summary CSV & Large-unseen SPO+ summary. \\
Step1c train-size 600 & \artifactpath{surrogate_experiment_results/Step1c/remote_results/formal_spoplus_ablation_val2000/train_size=600/metrics/spoplus_loss_curve.csv} & SPO+ training curve table. \\
Step1c train-size 600 & \artifactpath{surrogate_experiment_results/Step1c/remote_results/formal_spoplus_ablation_val2000/train_size=600/plots/2stage_mse_loss.png}; \artifactpath{plots/spoplus_loss.png} & Individual run curves. \\
Step1c train-size 1200 & \artifactpath{surrogate_experiment_results/Step1c/remote_results/formal_spoplus_ablation_val2000/train_size=1200/metrics/test_summary.csv} & Heldout400 SPO+ summary. \\
Step1c train-size 1200 & Internal large-unseen summary CSV & Large-unseen SPO+ summary. \\
Step1c train-size 1200 & \artifactpath{surrogate_experiment_results/Step1c/remote_results/formal_spoplus_ablation_val2000/train_size=1200/metrics/spoplus_loss_curve.csv} & SPO+ training curve table. \\
Step1c train-size 1200 & \artifactpath{surrogate_experiment_results/Step1c/remote_results/formal_spoplus_ablation_val2000/train_size=1200/plots/2stage_mse_loss.png}; \artifactpath{plots/spoplus_loss.png} & Individual run curves. \\
\end{longtable}

\FloatBarrier
\section{Main Conclusions}

\begin{enumerate}[label=\textbf{C\arabic*.}, leftmargin=2.2em]
    \item \textbf{The major new contribution is Step1c: a matched SPO+ surrogate
    ablation.}  The experiment changes the end-to-end surrogate from perturbed
    FY to reward-max SPO+ while keeping the Step1b protocol fixed as much as
    possible.

    \item \textbf{Validation2000 is now the strict checkpoint-selection
    protocol.}  This gives a fairer FY-vs-SPO+ comparison than using the older
    400-validation archive from Report0513.

    \item \textbf{On heldout400, SPO+ selected by validation SPO+ loss shows the
    strongest new signal.}  It improves over 2stage for all four Step1c train
    sizes and is competitive with, or better than, the FY validation-loss
    checkpoint in the available strict comparison.

    \item \textbf{On the large unseen test, the improvement is not yet robust.}  SPO+ is
    close to 2stage and FY, but the differences are small and should be
    described as near-tie or weak-effect rather than as a robust win.

    \item \textbf{The new diagnostics explain why this is not a simple story.}
    Many graphs have zero decision gap across methods; small average differences
    can be driven by the subset of graphs where the predicted weights change the
    selected KEP solution. This motivates the per-graph boxplots and nonzero-gap
    diagnostics.
\end{enumerate}

\section{Limitations and Next Steps}

\begin{enumerate}[leftmargin=2.2em]
    \item Complete and sync the missing Step1b validation2000 \(n=1200\) FY run,
    so the five-method comparison is available for all four train sizes.

    \item Add Step1c evaluation on the realistic-synthetic stress-test dataset.
    Current Step1c main figures focus on heldout400 and the large unseen test; the
    realistic2000 result is still a natural next panel.

    \item Repeat the strict Step1b/Step1c runs over multiple subset and theta
    seeds. The current result is a controlled single-seed comparison.

    \item Add a misspecification panel. The noisy-linear benchmark is relatively
    well specified for the two-feature probe, so a stronger decision-focused
    advantage may appear under utility-only probes, cPRA-only probes, nonlinear
    labels, or realistic-synthetic labels with the same two-feature model.

    \item Consider a warm-started SPO+ fine-tuning variant as an appendix
    experiment, but keep the main strict comparison with matched random
    initialization.
\end{enumerate}

\end{document}
